%----------------------------------------------------------------------------------------
% METADATA BLOCK (YAML)  — see inbox/README.md for every field
%----------------------------------------------------------------------------------------
% uid: ""
% school: "lycee"
% level: "4eme"
% section: "Mathématiques"
% chapter: "Coniques"
% title: "Une hyperbole qui passe par l'orthocentre"
% difficulty: "Difficile"
% tags: ["hyperbole-equilatere", "orthocentre", "cercle-d-euler", "coordonnees"]
% provenance: "adapted"
% batch: "2026-11-20-defis-geometrie-A"
%----------------------------------------------------------------------------------------

% ---STATEMENT---
Le plan est rapporté à un repère orthonormé $(\Omega;\vec{\imath},\vec{\jmath})$. On note $\Gamma$ l'hyperbole d'équation $xy=1$.
Soient $a$, $b$, $c$ trois réels non nuls deux à deux distincts et
$A\left(a,\dfrac1a\right)$, $B\left(b,\dfrac1b\right)$, $C\left(c,\dfrac1c\right)$ trois points de $\Gamma$.

% alt: Hyperbole d'équation xy=1, triangle ABC inscrit et orthocentre H sur l'hyperbole
\begin{center}
\begin{tikzpicture}[scale=0.85,line cap=round,line join=round,>=stealth]
\draw[->] (-3.2,0) -- (3.4,0) node[right]{$x$};\draw[->] (0,-3.2) -- (0,3.4) node[above]{$y$};
\draw[domain=0.3:3.3,smooth,samples=60,variable=\t] plot ({\t},{1/\t});
\draw[domain=-3.3:-0.3,smooth,samples=60,variable=\t] plot ({\t},{1/\t});
\draw (0.5,2) -- (2,0.5) -- (-1.5,-0.667) -- cycle;
\draw[dashed] (0.5,2) -- (0.667,1.5);
\fill(0.5,2) circle (1.4pt) node[above left] {$A$};
\fill(2,0.5) circle (1.4pt) node[above right] {$B$};
\fill(-1.5,-0.667) circle (1.4pt) node[below left] {$C$};
\fill(0.667,1.5) circle (1.4pt) node[above right] {$H$};
\node at (2.6,-0.5) {$\Gamma : xy=1$};
\end{tikzpicture}
\end{center}

\textbf{Partie A : l'orthocentre reste sur l'hyperbole.}
\begin{enumerate}
  \item Montrer que la droite $(BC)$ a pour coefficient directeur $-\dfrac{1}{bc}$. En déduire une équation de la hauteur du triangle $ABC$ issue de $A$.
  \item Montrer que le point $H\left(-\dfrac{1}{abc},\,-abc\right)$ est l'orthocentre de $ABC$.
  \item En déduire que $H\in\Gamma$.
\end{enumerate}

\textbf{Partie B : les centres des hyperboles équilatères passant par $A$, $B$, $C$.}

On oublie le repère de la partie A. Soit $ABC$ un triangle non rectangle, d'orthocentre $H$, inscrit dans un cercle $\mathcal C$ de centre $O$ et de rayon $R$. On note $N$ le milieu de $[OH]$.
On rappelle que le cercle d'Euler de $ABC$ a pour centre $N$ et pour rayon $\dfrac R2$.
Une hyperbole est dite \emph{équilatère} si ses asymptotes sont perpendiculaires.

\begin{enumerate}
  \item Soit $\mathcal H$ une hyperbole équilatère passant par $A$, $B$, $C$, de centre $\Omega$.
  \begin{enumerate}
    \item Justifier qu'on peut choisir un repère orthonormé d'origine $\Omega$ et une unité de longueur dans lesquels $\mathcal H$ a pour équation $xy=1$.
    \item Déduire de la partie A que $H\in\mathcal H$.
  \end{enumerate}
  \item On se place dans ce repère et on note $a$, $b$, $c$ les abscisses de $A$, $B$, $C$.
  \begin{enumerate}
    \item Montrer que $\mathcal C$ recoupe $\mathcal H$ en un quatrième point $D\left(\dfrac{1}{abc},\,abc\right)$, symétrique de $H$ par rapport à $\Omega$ (on pourra écrire que l'abscisse d'un point commun à $\mathcal C$ et $\mathcal H$ est racine d'un polynôme de degré $4$).
    \item En déduire que \textbf{le centre $\Omega$ de toute hyperbole équilatère passant par $A$, $B$, $C$ appartient au cercle d'Euler de $ABC$.}
  \end{enumerate}
\end{enumerate}

% ---SOLUTION---
\textbf{Partie A.}
\begin{enumerate}[itemsep=14pt]
  \item\leavevmode
  \begin{prop}[Coefficient directeur et perpendicularité]
  La droite passant par $M(x_M,y_M)$ et $N(x_N,y_N)$, avec $x_M\neq x_N$, a pour coefficient directeur \ib{$\dfrac{y_N-y_M}{x_N-x_M}$}.

  Deux droites de coefficients directeurs $m$ et $m'$ sont perpendiculaires si et seulement si \ib{$mm'=-1$}.
  \end{prop}
  Coefficient directeur de $(BC)$ :
  \[\dfrac{\dfrac1c-\dfrac1b}{c-b}=\dfrac{\dfrac{b-c}{bc}}{c-b}=-\dfrac{1}{bc}.\]

  La hauteur issue de $A$ est perpendiculaire à $(BC)$ : son coefficient directeur $m$ vérifie $m\times\left(-\dfrac1{bc}\right)=-1$, donc $m=bc$.

  Elle passe par $A\left(a,\dfrac1a\right)$ : \ib{$y-\dfrac1a=bc\,(x-a)$}.

  \item On remplace $H\left(-\dfrac1{abc},-abc\right)$ dans l'équation de la hauteur issue de $A$.

  Membre de gauche : $y_H-\dfrac1a=-abc-\dfrac1a$.

  Membre de droite : $bc\,(x_H-a)=bc\left(-\dfrac1{abc}-a\right)=-\dfrac1a-abc$.

  Les deux membres sont égaux : $H$ est sur la hauteur issue de $A$.

  Les coordonnées de $H$ sont \emph{symétriques} en $a$, $b$, $c$. En échangeant les rôles de $A$, $B$, $C$, le même calcul montre que $H$ est aussi sur la hauteur issue de $B$ et sur celle issue de $C$.

  $H$ est donc le point commun aux trois hauteurs : \ib{$H$ est l'orthocentre de $ABC$}.

  \item $x_Hy_H=\left(-\dfrac{1}{abc}\right)(-abc)=1$, donc \ib{$H\in\Gamma$}.
\end{enumerate}

\textbf{Partie B.}
\begin{enumerate}[itemsep=14pt]
  \item\leavevmode
  \begin{prop}[Hyperbole rapportée à ses asymptotes]
  Si les asymptotes d'une hyperbole de centre $\Omega$ sont les axes d'un repère orthonormé d'origine $\Omega$, son équation est \ib{$xy=k$} avec $k\neq0$.
  \end{prop}
  \begin{enumerate}[itemsep=8pt]
    \item $\mathcal H$ est équilatère : ses asymptotes sont deux droites perpendiculaires qui se coupent en $\Omega$. On les prend pour axes.

    Alors $\mathcal H$ a pour équation $xy=k$, $k\neq0$.

    Si $k<0$, on change l'orientation de l'axe des ordonnées ($y$ devient $-y$) : on obtient $k>0$.

    On prend alors $\sqrt k$ pour unité de longueur : $X=\dfrac{x}{\sqrt k}$, $Y=\dfrac{y}{\sqrt k}$ donnent $XY=1$.
    \item Dans ce repère, $A$, $B$, $C$ sont trois points de $\Gamma:xy=1$. L'orthocentre est une notion géométrique : il ne dépend ni du repère ni de l'unité. Le résultat de la partie A donne $H\in\mathcal H$.
  \end{enumerate}

  \item\leavevmode
  \begin{prop}[Relations entre coefficients et racines]
  Si un polynôme unitaire de degré $4$ s'écrit $x^4+\cdots+e$ et a pour racines $x_1,x_2,x_3,x_4$, alors \ib{$x_1x_2x_3x_4=e$} (terme constant).
  \end{prop}
  \begin{enumerate}[itemsep=8pt]
    \item Le cercle $\mathcal C$ passe par $A$, $B$, $C$ : il a une équation $x^2+y^2+px+qy+s=0$.

    Un point de $\mathcal H$ a pour ordonnée $\dfrac1x$ ($x\neq0$). Il est sur $\mathcal C$ si et seulement si
    \[x^2+\dfrac1{x^2}+px+\dfrac qx+s=0.\]
    On multiplie par $x^2$ : \ib{$x^4+px^3+sx^2+qx+1=0$}\quad $(E)$.

    $A$, $B$, $C$ sont communs à $\mathcal C$ et $\mathcal H$ : $a$, $b$, $c$ sont trois racines distinctes de $(E)$. Le polynôme du premier membre est unitaire de degré $4$ ; sa quatrième racine $d$ est réelle et
    \[(E)\iff (x-a)(x-b)(x-c)(x-d)=0.\]
    Le terme constant vaut $1$ : $abcd=1$, donc $d=\dfrac1{abc}$.

    Le point $D\left(d,\dfrac1d\right)=\left(\dfrac1{abc},abc\right)$ est sur $\mathcal H$ et son abscisse est racine de $(E)$ : $D\in\mathcal C$.

    Enfin $D=(-x_H,-y_H)$ : $D$ est le symétrique de $H$ par rapport à $\Omega$.
    \item $N$ est le milieu de $[OH]$ : $\overrightarrow{\Omega N}=\dfrac12\left(\overrightarrow{\Omega O}+\overrightarrow{\Omega H}\right)$.

    $\Omega$ est le milieu de $[HD]$ : $\overrightarrow{\Omega D}=-\overrightarrow{\Omega H}$.

    On remplace $\overrightarrow{\Omega H}$ par $-\overrightarrow{\Omega D}$ : \textbf{$\overrightarrow{\Omega N}=\dfrac12\left(\overrightarrow{\Omega O}-\overrightarrow{\Omega D}\right)=\dfrac12\overrightarrow{DO}$}.

    Comme $D\in\mathcal C$, $DO=R$ : \ib{$\Omega N=\dfrac R2$}.

    $\Omega$ est sur le cercle de centre $N$ et de rayon $\dfrac R2$, c'est-à-dire sur le cercle d'Euler. $\blacksquare$
  \end{enumerate}
\end{enumerate}

\textbf{Idée :} dans le repère centré sur l'hyperbole, le quatrième point commun avec le cercle circonscrit est le symétrique de $H$. Alors $\Omega$ est milieu de $[HD]$, $N$ milieu de $[HO]$ : la droite des milieux donne $\Omega N=\dfrac{OD}2=\dfrac R2$.

\textbf{Piège :} les racines de $(E)$ sont comptées avec multiplicité : si $\mathcal C$ est tangent à $\mathcal H$ en un sommet, $d$ est égal à l'un des réels $a$, $b$, $c$ ; la formule $d=\dfrac1{abc}$ reste valable.
